\subsection{局部环}

命题 1. 设 A 是环，I 是 A 的不可逆元素组成的集合。集合 I 是 A 的异于 A 的理想的并。此外，下列条件等价：

\begin{enumerate}
\def\labelenumi{(\alph{enumi})}
\item
  I 是理想。
\item
  \(\mathbf{A}\) 的异于 \(\mathbf{A}\) 的理想组成的集合有一个最大元素。
\item
  A 有唯一的极大理想。
\end{enumerate}

关系 \(x \in I\) 等价于 \(1 \notin xA\)，从而 \(xA \neq A\)。若 \(a\) 是 \(A\) 的异于 \(A\) 的理想且 \(x \in a\)，则 \(xA \subset a\)，从而 \(xA \neq A\) 且 \(x \in I\)。因而 \(A\) 的每个异于 \(A\) 的理想都包含于 \(I\)，且每个元素 \(x \in I\) 都属于某个主理想 \(xA \neq A\)。这就证明了第一个断言，它立即蕴含 (a)、(b) 与 (c) 的等价性。

注 (1). 注意，若 (c) 成立，则 I 是环 A 的 Jacobson 根（《代数学》第 VIII 章，§ 6，第 3 节，定义 3）。

定义 1. 环 A 称为局部环，如果它满足命题 1 的等价条件 (a)、(b) 与 (c)。A 对其 Jacobson 根（它就是 A 的唯一极大理想）的商称为 A 的剩余域。

定义 2. 设 A、B 是两个局部环，m、n 是它们各自的极大理想。同态 \(u: \mathbf{A} \to \mathbf{B}\) 称为局部的，如果 \(u(\mathfrak{m}) \subset \mathfrak{n}\)。

这等价于说 \(\bar{u}^{1}(\mathfrak{n}) = \mathfrak{m}\)，因为 \(\bar{u}^{1}(\mathfrak{n})\) 此时是一个包含 m 而不包含 1 的理想，从而等于 m。取商，我们于是由 u 典范地导出一个从 A 的剩余域到 B 的剩余域的单同态 A/m → B/n。

\textbf{例}\par

\begin{enumerate}
\def\labelenumi{(\arabic{enumi})}
\item
  域是局部环。化为 0 的环不是局部环。
\item
  设 A 是局部环，k 是它的剩余域。形式幂级数环 \(B = A[[X_{1}, \ldots, X_{n}]]\) 是局部环，因为 B 的不可逆元素恰是常数项在 A 中不可逆的形式幂级数（《代数学》第 IV 章，§ 5，第 6 节，命题 4）。A 到 B 的典范单射是局部同态，且相应的剩余域的单射是同构。
\item
  设 \(\mathfrak{b}\) 是环 A 的理想，它只包含于单独一个极大理想 m 中；则 A/b 是局部环，其极大理想是 m/b，剩余域典范地同构于 A/m。这特别适用于 \(b = m^{k}\) 的情形，其中 m 是 A 的任意极大理想（\(\S\) 1，第 1 节，命题 1 的推论）。若 A 本身是以 m 为极大理想的局部环，则对 A 的每个理想 \(b \neq A\)，A/b 是局部环，典范同态 \(A \to A/b\) 是局部同态，且相应的剩余域同态是同构。
\item
  设 X 是拓扑空间，\(x_{0}\) 是 X 的一点，A 是在点 \(x_{0}\) 处的芽组成的环，这些芽来自在 \(x_{0}\) 的某个邻域中连续的实值函数（《一般拓扑学》第 I 章，§6，第 10 节）。显然，为使连续函数 f 在 \(x_{0}\) 处的芽在 A 中可逆，必须且只需 \(f(x_{0}) \neq 0\)，因为这蕴含 \(f(x) \neq 0\) 在 \(x_{0}\) 的某个邻域中成立。因此环 A 是局部环，其极大理想 m 是在 \(x_{0}\) 处取零值的函数的芽组成的集合；取商，A 到 R 的映射 \(g \mapsto g(x_{0})\) 给出剩余域 A/m 到 R 的一个同构。
\end{enumerate}

命题 2. 设 A 是环，p 是 A 的素理想。环 \(A_{p}\) 是局部的；其极大理想是理想 \(pA_{p} = p_{p}\)（由 p 在 \(A_{p}\) 中的典范像生成）；其剩余域典范地同构于 A/p 的分式域。

设 \(S = A - p\)，并设 \(j: A \to A_p\) 是典范同态；\(p\) 是素理想的假设蕴含 \(p\) 关于 \(S\) 是饱和的，因而 \(j^{1}(\mathfrak{pA}_{\mathfrak{p}}) = \mathfrak{p}\)（§ 2，第 4 节，命题 10），而且由于 A 的不与 S 相交的理想恰是包含于 \(\mathfrak{p}\) 中的那些，前两个断言是 § 2 第 5 节命题 11 (ii) 的特殊情形。此外，若 \(f\) 是典范同态 \(A \to A/\mathfrak{p}\)，则 \(f(S)\) 是 \(\neq 0\) 的元素组成的集合（整环 \(A/\mathfrak{p}\) 的），因而最后一个断言是 § 2 第 5 节命题 11 (i) 的特殊情形。

定义 3. 设 A 是环，p 是 A 的素理想。环 \(A_{p}\) 称为 A 在 p 处的局部环，或在无歧义时称为 p 的局部环。

注 (2). 若 A 是局部环而 m 是其极大理想，则 A --- m 的元素都是可逆的（命题 1），因而 \(A_{m}\) 典范地等同于 A（§ 2，第 1 节，注 5）。

\textbf{例}\par

\begin{enumerate}
\def\labelenumi{(\arabic{enumi})}
\setcounter{enumi}{4}
\tightlist
\item
  设 \(p\) 是素数。局部环 \(\mathbf{Z}_{(p)}\) 是有理数 \(a / b\) 组成的集合，其中 \(a, b\) 是有理整数且 \(b\) 与 \(p\) 互素；\(\mathbf{Z}_{(p)}\) 的剩余域同构于素域 \(\mathbf{F}_p = \mathbf{Z} / (p)\)。
\end{enumerate}

* (6) 设 V 是仿射代数簇，A 是 V 上的正则函数环，W 是 V 的不可约子簇，而 p 是 A 中由在 W 的每点取零值的函数组成的（必然为素的）理想。环 \(\mathbf{A}_{\mathfrak{p}}\) 称为 V 上 W 的局部环。*

命题 3. 设 A 是环，p 是 A 的素理想，且 S = A - p.~对每个理想 \(\mathfrak{b}'\)（\(\mathbf{A}_{\mathfrak{p}}\) 的，异于 \(\mathbf{A}_{\mathfrak{p}}\)），设 \(\mathfrak{b}\) 是 A 的理想 \((i_{\mathbf{A}}^{\mathbf{S}})^{-1}(\mathfrak{b}')\)，于是 \(\mathfrak{b}' = \mathfrak{bA}_{\mathfrak{p}}\)。

\begin{enumerate}
\def\labelenumi{(\roman{enumi})}
\item
  设 \(f\) 是典范同态 \(\mathbf{A} \to \mathbf{A} / \mathfrak{b}\)。从 \(\mathbf{A}_{\mathfrak{p}}\) 到 \((\mathbf{A} / \mathfrak{b})_{\mathfrak{p} / \mathfrak{b}}\) 的与 \(f\) 典范相联的同态（§ 2，第 1 节，命题 2）是满射的，其核是 \(\mathfrak{b}'\)，由此通过取商定义了 \(\mathbf{A}_{\mathfrak{p}} / \mathfrak{b}'\) 到 \((\mathbf{A} / \mathfrak{b})_{\mathfrak{p} / \mathfrak{b}}\) 的一个典范同构。
\item
  映射 \(\mathfrak{b}' \to \mathfrak{b} = (i_{\mathbf{A}}^{\mathbb{S}})^{-1}(\mathfrak{b}')\) 限制到 \(\mathbf{A}_{\mathfrak{p}}\) 的素理想集合上，是这个集合到 \(\mathbf{A}\) 的素理想集合（包含于 \(\mathfrak{p}\) 中的那些）的（关于包含关系的）同构。若 \(\mathfrak{b}'\) 在 \(\mathbf{A}_{\mathfrak{p}}\) 中是素的，则存在环 \(\mathbf{A}_{\mathfrak{b}}\) 到环 \((\mathbf{A}_{\mathfrak{p}})_{\mathfrak{b}'}\) 的同构，它把 \(a / s\) 映为 \((a / 1) / (s / 1)\)（对所有 \(a \in \mathbf{A}\)、\(s \in \mathbf{A} - \mathfrak{b}\)）。
\end{enumerate}

这只是 § 2 第 5 节命题 11 的一个特殊情形。

\textbf{注}\par

\begin{enumerate}
\def\labelenumi{(\arabic{enumi})}
\setcounter{enumi}{2}
\item
  若 \(\mathfrak{a}\) 是 \(\mathbf{A}\) 的不包含于 \(\mathfrak{p}\) 的理想，则 \(\mathfrak{a}\mathrm{A}_{\mathfrak{p}} = \mathrm{A}_{\mathfrak{p}}\) 且 \((\mathrm{A} / \mathfrak{a})_{\mathfrak{p}} = 0\)（§ 2，第 5 节，注）。
\item
  设 A、B 是两个环，\(\rho: A \to B\) 是同态，\(q\) 是 B 的素理想，而 \(p\) 是 A 的素理想 \(\rho^{-1}(q)\)。由于 \(\rho(A - p) \subset B - q\)，一个典范同态 \(\rho_q: A_p \to B_q\) 由 \(\rho\) 导出（\(\S 2\)，第 1 节，命题 2），且立即可见 \(\rho_q(pA_p) \subset qB_q\)，因而 \(\rho_q\) 是局部同态。
\end{enumerate}
% semantic-fragments: legacy-input-seam
\relax{}
